% ============================================================================
%  Chapter 1. Introduction
%
%  \input file. No preamble. The parent document must load:
%      \usepackage{booktabs}   % \toprule \midrule \bottomrule \cmidrule
%      \usepackage{amsmath}    % aligned display math
%      \usepackage{graphicx}   % (only if figure assets are substituted in)
%      \usepackage[round]{natbib}  % \citep \citet
%  Bibliography: accord/thesis/refs.bib
%
%  Every quantity in this chapter is forward-referenced to the chapter that
%  derives it. No number appears here that is not derived in Chapters 3 to 8 or
%  recorded in accord/CORRECTIONS.md.
% ============================================================================

\chapter{Introduction}
\label{ch:intro}

\section{A Patient Goes Home on a Held Medication}
\label{sec:intro-problem}

A patient admitted for decompensated heart failure is discharged on eleven
medications. Somewhere in the middle of the stay a nurse records that the loop
diuretic was withheld because the patient had become hypovolemic. Two days
later the discharge summary, written by a physician who is reconciling a
medication list against an admission list, instructs that the same diuretic be
continued at its admission dose. The two documents disagree. Only one of them
governs what the patient does after leaving the building, and it is not the one
that contains the reason the instruction is wrong.

Nothing about this scenario is exotic. Polypharmacy is the norm in heart
failure, the drugs involved have narrow therapeutic windows, and the transition
from inpatient to outpatient care is the single most error-prone moment in the
medication lifecycle. Unintended medication discrepancies are present in a
large share of hospital admissions \citep{cornish2005discrepancies} and
discharges \citep{wong2008reconciliation}; adverse events in the weeks after
discharge are common, frequently drug-related, and frequently preventable
\citep{forster2003adverse}; errors traceable to discontinuity between the
inpatient and outpatient settings affect a substantial fraction of discharged
patients \citep{moore2003discontinuity}; and the information transfer that
would prevent them is unreliable in both content and timing
\citep{kripalani2007deficits}. Adverse drug events were established as a
first-order inpatient safety problem three decades ago
\citep{bates1995ade} and the transition-of-care component of that problem has
not been solved since.

This thesis takes as its object one narrow, checkable instance of that problem:
a \emph{cross-note medication contradiction}, defined as a disagreement about a
specific drug between a nursing progress note and the discharge summary of the
same admission. Chapter~\ref{ch:candidates} constructs four transparent rules
over extracted medication mentions and applies them to MIMIC-III
\citep{johnson2016mimic}, yielding $1{,}978$ candidate contradictions across
$1{,}203$ admissions drawn from a heart-failure cohort of $9{,}783$. A
board-certified pharmacist annotating a stratified sample of those candidates
confirmed $12$ genuine contradictions in $299$ pairs, four of them rated $4$ or
$5$ on a five-point danger scale.

Those numbers establish the tension this thesis is built on. The phenomenon is
systematic rather than anecdotal: rules over a single institution's notes
surface nearly two thousand candidates, and a domain expert confirms that some
of them are real and dangerous. But the phenomenon is also \emph{rare} in the
population a system would have to screen, and rarity is not a nuisance
parameter here. It is the thing that breaks every stage of the pipeline that
follows.

\section{The Pipeline the Field Would Build, and Where It Breaks}
\label{sec:intro-gap}

Faced with this problem, the standard construction is close to automatic. Mine
candidates with cheap, high-recall rules, because exhaustive expert review of
every cross-note medication pair is infeasible. Label the candidates with an
instruction-tuned language model, because expert annotation does not scale and
recent work reports that such models match or exceed crowd annotators on many
classification tasks \citep{gilardi2023chatgpt}. Distil those labels into a
purpose-built neural detector over a domain-pretrained clinical encoder
\citep{gu2021biomedbert,alsentzer2019clinicalbert}, with architectural
components encoding the domain structure that generic natural language
inference models lack. Evaluate against the expert annotations. Report.

The construction is defensible at every step, and this thesis executed it. It
fails at every step, and the failures are not the ones the literature would
predict.

\paragraph{The gap.} Prior work on clinical natural language inference
establishes that encoder models can classify entailment relations between
clinician-authored sentence pairs \citep{romanov2018mednli}, and prior work on
weak supervision establishes that noisy programmatic labels can substitute for
hand annotation when the noise sources are multiple and their errors are
conditionally independent \citep{ratner2016dataprogramming,ratner2017snorkel}.
Neither body of work addresses what happens when the positive class occupies
four percent of an already-enriched candidate pool, when there is exactly one
noise source, and when that source's error is a systematic prior shift rather
than independent noise. Under those conditions the standard recipe does not
degrade gracefully. It produces a system whose outputs are statistically
independent of the labels, and it produces that system while emitting metrics
that look survivable.

\paragraph{The consequence for what gets reported.} The reason this matters
beyond one project is that the failure is quiet. A system with no signal, trained
on supervision whose marginal is inverted, still produces a contradiction
$F_1$, a recall figure, an ablation table, and an auxiliary-task accuracy. Every
one of those numbers can be reported in good faith by someone who has not run
the specific diagnostic that would falsify it. Chapters~\ref{ch:arch}
and~\ref{ch:eval} document four such claims from this project's own earlier
submission, each of which survived internal review and none of which survives
its own first check.

\section{Research Questions}
\label{sec:intro-rqs}

The project was designed around four questions. They are stated here in the
form in which they were posed, before the answers were known, because the value
of the thesis lies substantially in the gap between what was asked and what
turned out to be answerable.

\begin{description}
\item[RQ1 (Authority).] Does conditioning a cross-note inference model on
      \emph{note-type authority} outperform symmetric natural language
      inference? The construct, developed in Chapter~\ref{ch:data} as
      \emph{clinical note priority}, is that a discharge summary and a nursing
      progress note carry different institutional weight, so a disagreement
      between them is not a symmetric relation.
\item[RQ2 (Temporality).] Does encoding temporal metadata, the ordering of the
      two notes and the interval between them, improve detection over a
      time-agnostic pairwise model?
\item[RQ3 (Severity).] Can published clinical practice guidelines be used to
      rank confirmed contradictions by danger, so that a triage system surfaces
      the harmful ones first?
\item[RQ4 (Supervision).] Can weak supervision from an instruction-tuned
      language model replace expert annotation at the scale required to train a
      detector for a rare clinical event?
\end{description}

\section{Contributions, Stated as Findings}
\label{sec:intro-contributions}

This thesis contributes no working system. Its contributions are measurements,
and they are stated as such. Each is a finding about the pipeline rather than a
component of it, each is traceable to a named chapter, and each is accompanied
by the diagnostic that would have produced it earlier.

\begin{enumerate}
\item \textbf{Rule-based candidate mining yields a four percent precision pool,
      and the evidence attached to those candidates was frequently the wrong
      evidence.} Four transparent rules enrich the candidate space to a
      pharmacist-measured base rate of $4.0\%$, which is a real and useful
      enrichment over an unfiltered background rate orders of magnitude lower.
      But the mechanism that decided \emph{which} candidate a sentence pair was
      flagged for was independent of the mechanism that selected the sentence
      pair to display, and for $41.4\%$ of action-conflict candidates the
      displayed evidence did not contain the flagged conflict. Every downstream
      consumer, the annotator, the weak labeller, and the model, reasoned about
      that evidence. The $4.0\%$ figure is therefore an upper bound on the
      validity of the annotation exercise rather than a clean estimate of pool
      precision. Chapter~\ref{ch:candidates}.

\item \textbf{Language model weak supervision is miscalibrated by a factor of
      roughly twenty-one, and the detector inherits the labeller's prior rather
      than its judgement.} The weak labeller assigned the positive class to
      $84.9\%$ of $1{,}679$ candidates against a measured base rate of $4.0\%$.
      Under that base rate at least $95\%$ of its positive labels are wrong. The
      trained detector then reproduced the labeller's per-rule
      predicted-positive rate to within $1.5$ percentage points on three of four
      rules, while the rules' true prevalences differ from one another by more
      than a factor of four. Reproducing the quantity that carries no
      information about the truth, and not the quantity that does, is the
      signature of a model that has learned the label's marginal rather than its
      conditional. This answers RQ4 in the negative, and the negative is
      specific: the failure is calibration, not comprehension, because the same
      labeller recovers $11$ of the $12$ true contradictions.
      Chapter~\ref{ch:weak}.

\item \textbf{Authority conditioning contributes nothing, and the claim that it
      does is falsified four independent ways.} The authority parameter reported
      as \emph{discovered} at $(\alpha_N, \alpha_D) \approx (0.5, 1.0)$ was
      \emph{initialised} at exactly $[0.5, 1.0]$; it moved by $1.8\%$ and
      $0.6\%$, agreeing to the third decimal across seeds; it sits on a branch
      that a subsequent LayerNorm largely erases, leaving a near-zero gradient;
      and the project's own ablation shows that constraining
      $\alpha_N = \alpha_D$ \emph{improves} both reported metrics. RQ1 is
      therefore not answered in the affirmative and, more precisely, was never
      tested: the implementation could not have detected an authority signal had
      one existed. Chapter~\ref{ch:arch}.

\item \textbf{The trained detector carries no rank signal, and the finding
      survives the one defect that could plausibly have manufactured it.} On the
      $299$ annotated pairs the contradiction probability achieves AUROC
      $0.460$, average precision at the no-skill line, Mann-Whitney $p = 0.64$,
      and a bootstrap $95\%$ confidence interval of $[0.287, 0.633]$. The
      fifteen highest-scored pairs contain zero true contradictions. Restricted
      to Rules 3 and 4, the strata that are free of the exemplar defect by
      construction, AUROC is $0.5078$ with Mann-Whitney $p = 0.953$, which is
      chance to four significant figures. That stratum contains six positives,
      so the estimate is wide, and it is reported with that width attached.
      Chapters~\ref{ch:eval} and~\ref{ch:discussion}.

\item \textbf{RQ3 is untestable from this work, because the component that would
      have answered it was never trained.} The guideline-grounded harm estimator
      has no guideline corpus; its ``guideline bank'' is twenty randomly
      initialised vectors. Its harm head appears in no term of the training
      objective. Its outputs across the evaluation set span $0.398$ to $0.610$,
      the range of an untrained sigmoid, and correlate with the pharmacist's
      danger rating at $\rho = -0.0001$. That null is a measurement of an
      untrained module, not evidence about whether clinical severity is
      predictable, and reporting it as the latter would be the single most
      misleading move available in this document. Chapters~\ref{ch:arch}
      and~\ref{ch:eval}.

\item \textbf{A prescriptive result about evaluation under low prevalence.} At
      the evaluation set's operating prevalence of $12$ positives against $241$
      negatives, balanced accuracy, Matthews correlation, and
      prevalence-adjusted bias-adjusted kappa all prefer a system that recovers
      $6$ of $12$ contradictions over one that recovers $11$, and the exchange
      rate is available in closed form: the preference flips once the additional
      true negatives per forgone true positive exceed $Q/P = 20.08$.
      Chapter~\ref{ch:metrics} derives this, and then sharpens it against the
      empirical situation, in which the systems being compared mostly have no
      discriminative ability at all.
\end{enumerate}

\paragraph{RQ2 has no answer, for a reason that is itself a finding.}
\texttt{CHARTTIME} is null for $100\%$ of MIMIC-III discharge summaries. The
temporal ordering RQ2 requires cannot be established at finer than day
granularity between the two notes of a pair, and the filter that purported to
measure whether conflicts were resolved by intervening documentation was, on
inspection, matching a twenty-way alternation containing both directions of the
medication decision at once. The temporal question is not answered negatively.
It is unaddressable in this resource, and the correct report is the data
limitation rather than the measurement. Chapter~\ref{ch:data},
Section~\ref{sec:charttime}.

\section{What the Results Look Like}
\label{sec:intro-results}

The headline is stated once, plainly, so that nothing later in the document
reads as a reveal.

On the $299$ pharmacist-annotated pairs, the trained detector predicts
\emph{contradiction} for $278$, achieving precision $0.047$ against a base rate
of $0.040$, recall $0.917$, and accuracy $0.115$. At thresholds of $0.05$,
$0.10$, and $0.15$ it predicts positive for every pair, because no score in the
evaluation set falls below $0.15$. Its Matthews correlation is $-0.007$. The
weak labeller it was distilled from scores $-0.044$, which is worse than
chance. A general-purpose model prompted zero-shot on the same $299$ pairs, with
no training data and no clinical components, is the only system in the
comparison with a positive Matthews correlation, at $+0.109$ and $8.45\%$
precision.

Two consequences follow immediately and are developed at length later. First,
the task is not impossible: something above chance is extractable from the same
evidence, so the failure is a property of this pipeline rather than of the
problem. Second, the failures are causally independent of one another. The
exemplar defect, the prevalence mismatch, the initialisation artifact, the
auxiliary-head leakage, and the untrained severity head each have their own
mechanism and their own detector, and repairing any one of them in isolation
would still leave a system with AUROC near $0.5$.

A third consequence is the one this thesis regards as most transferable. The
component sanity check in Chapter~\ref{ch:eval} shows that BioMedBERT
fine-tuned on MedNLI reaches $83.90\%$ accuracy and $83.91\%$ macro $F_1$,
inside the published range for comparable encoder systems. That result is a
property of the backbone and contains zero components of the system under
study. Reporting it under the system's name would have produced a document in
which a strong number and a failing system coexist without the reader being
able to tell which is which. This thesis does not do that, and
Chapter~\ref{ch:eval} explains why the temptation to do it is structural rather
than personal.

\section{Scope, and What This Thesis Does Not Claim}
\label{sec:intro-scope}

Three boundaries are worth fixing before the reader encounters the chapters that
test them.

\paragraph{This is a negative result about one pipeline on one resource, not
about the problem.} Nothing here establishes that cross-note contradiction
detection is unsolvable, that clinical note priority is unlearnable, or that
clinical severity is unpredictable. Each of those hypotheses remains open, and
Chapter~\ref{ch:discussion} states what a correct test of each would require.
What is established is that this pipeline did not test them, and that the
measurements which would have revealed that were cheap and were not run.

\paragraph{The evidence base is small and single-site.} Twelve positives
support every interval in Chapter~\ref{ch:eval}. One annotator produced every
gold label, with a two-pass protocol and a two-week washout, and inter-rater
reliability is not established. A previously reported $\kappa = 1.000$ is
retracted as circular. The cohort is critical-care admissions at one US academic
medical centre over a historical period. Chapter~\ref{ch:discussion} treats
these as first-order limitations rather than as caveats.

\paragraph{The system as measured is not deployable, and its failure mode is
worse than inaction.} At $91.7\%$ recall and $4.7\%$ precision, a clinical
alert built on this detector generates roughly twenty false alarms per true
contradiction. Alert override rates in computerised order entry are already
between $49\%$ and $96\%$ \citep{vandersijs2006overriding}, and alerting burden
degrades response \citep{ancker2017alertfatigue}. A system that consumes
clinician attention while carrying, by Chapter~\ref{ch:eval}'s Measurement 1,
no information, is not a weak system. It is a harmful one.

\section{Structure of the Thesis}
\label{sec:intro-structure}

Chapter~\ref{ch:background} organises the prior work into five clusters,
identifies the six gaps the project set out to address, and is candid that two
of those gaps are weakened as framings by this thesis's own findings. It also
develops three bodies of literature that the project needed and did not consult:
class imbalance, reporting standards for clinical prediction models, and the
behaviour of agreement and correlation metrics at low prevalence.

Chapters~\ref{ch:data} through~\ref{ch:eval} are the empirical core.
Chapter~\ref{ch:data} builds the cohort and the medication mention table and
documents two constraints, the \texttt{CHARTTIME} absence and an $8.9{:}1$
asymmetry in extraction yield between the two note types, that govern everything
after it. Chapter~\ref{ch:candidates} constructs the four rules, reports the
gold standard, and measures the exemplar selection failure.
Chapter~\ref{ch:weak} reports the weak supervision prevalence mismatch and the
transfer of the labeller's marginal to the student. Chapter~\ref{ch:arch}
audits the model against the code that produced its results and falsifies three
claimed contributions. Chapter~\ref{ch:eval} measures the system four
independent ways and reports that all four agree.

Chapter~\ref{ch:metrics} takes up the evaluation question separately, because it
generalises past this system: it derives in closed form what balanced accuracy,
Matthews correlation, and prevalence-adjusted kappa are actually selecting for
at this prevalence, and then sharpens the argument against the empirical
situation in which most of the compared systems have no signal to select among.
Chapter~\ref{ch:discussion} states the limitations without hedging, including
one design flaw serious enough to invalidate the intended deployment story, and
sets out what a continuation of this work should do first.
Chapter~\ref{ch:conclusion} is short.

\section{Proposed Figures}
\label{sec:intro-figures}

\begin{figure}[t]
\centering
\framebox[\linewidth]{\parbox[c][6cm][c]{0.92\linewidth}{\centering
\textit{[Figure asset 1.1: the pipeline as a left-to-right chain with a status}
\textit{badge under each stage. Stages: cohort construction; medication}
\textit{extraction; rule-based candidate generation; LLM weak supervision;}
\textit{detector training; evaluation. Badges, in order: WORKS; UNVALIDATED;}
\textit{ENRICHES BUT MIS-ATTRIBUTES EVIDENCE (41.4\%); FAILS (21x prevalence}
\textit{mismatch); THREE CLAIMS FALSIFIED; FAILS (AUROC 0.460). Beneath each}
\textit{badge, in small type, the single diagnostic that would have caught the}
\textit{failure at that stage. No clinical text.]}}}
\caption{The thesis in one view. The reader should notice that the failures are
at different stages, have different mechanisms, and each have a cheap detector,
and that a project fixing any one of them in isolation would still have produced
a system with no ranking signal.}
\label{fig:pipeline-status}
\end{figure}

\begin{figure}[t]
\centering
\framebox[\linewidth]{\parbox[c][5cm][c]{0.92\linewidth}{\centering
\textit{[Figure asset 1.2: three horizontal bars on a shared 0--100\% axis.}
\textit{Bar 1, weak labeller positive rate: 84.9\%. Bar 2, detector}
\textit{predicted-positive rate on the gold set: 278/299. Bar 3,}
\textit{pharmacist-measured true rate: 4.0\%. A bracket spans bars 1 and 3}
\textit{annotated ``21x''. Inset: MCC for four systems, GPT-4o zero-shot}
\textit{+0.109, GPT-4o few-shot +0.012, ACCORD -0.007, Llama 3 8B -0.044,}
\textit{with a zero line marked ``chance''.]}}}
\caption{The single number that determines the rest of the thesis, and its
consequence. A labeller whose marginal is inverted relative to the truth
produces a student that reproduces the marginal, and a system whose Matthews
correlation is at or below chance while a zero-shot general-purpose model on the
same evidence is above it.}
\label{fig:intro-headline}
\end{figure}

\section{Summary}
\label{sec:intro-summary}

Dangerous medication contradictions between nursing notes and discharge
summaries are a real and recurring patient-safety problem. Four transparent
rules over one institution's notes surface $1{,}978$ candidates across $1{,}203$
admissions, and a board-certified pharmacist confirms $12$ genuine
contradictions in a $299$-pair sample, four of them rated $4$ or $5$ for danger.

The pipeline the field would naturally build for this problem, rule-based
candidate mining followed by language model weak supervision followed by a
purpose-built neural detector, fails at every stage. The rules attach the wrong
evidence to $41.4\%$ of action-conflict candidates. The weak labeller is
positive at $84.9\%$ against a $4.0\%$ base rate and the detector inherits that
marginal rather than any judgement. Three architectural contributions do not
survive an audit of the code that produced their numbers, and a fourth research
question was never tested because the component that would answer it received no
gradient. The resulting detector has AUROC $0.460$ on the full evaluation set
and $0.5078$ on the strata that are free of the evidence-selection defect by
construction.

These failures are measurable, they are diagnosable, and each of them was
detectable before any results table existed. That is the claim of this thesis,
and Chapters~\ref{ch:data} through~\ref{ch:metrics} are its evidence.
